\section*{Capítulo \ref{ch:b3:complete} --- Espacios completos:
Baire, Ascoli, Stone--Weierstrass}

\begin{solution}{exo:b3:complete:1}
(a) Para $m \geq n$, $f_n - f_m$ se anula fuera de un intervalo de
longitud $\frac1n$ y está acotada por $1$: $\norm{f_n - f_m}_1
\leq \frac1n \to 0$: de Cauchy. Si $f_n \to f$ en $\norm\cdot_1$ con $f$
\omterm{def:b3:topology:continuity}{continua}: para $\alpha <
\frac12$ fijo, $\int_0^{\alpha}\abs f = \lim \int_0^\alpha\abs{f -
f_n + f_n} \leq \lim\bigl(\norm{f - f_n}_1 + 0\bigr) = 0$
($f_n \equiv 0$ allí para $n$ grande), luego $f \equiv 0$ sobre
$[0, \frac12)$ (\omterm{def:b3:topology:continuity}{continuidad});
análogamente $f \equiv 1$ sobre $(\frac12, 1]$: ninguna función
\omterm{def:b3:topology:continuity}{continua} cumple ambas cosas. Así
que el espacio es incompleto --- el
\omterm{thm:b3:complete:completion}{completado} es $L^1$, construido en
\cref{ch:b3:lp}.

(b) Sea $(x_n)$ de Cauchy; elíjanse $n_1 < n_2 < \cdots$ con
$\norm{x_{n_{k+1}} - x_{n_k}} \leq 2^{-k}$. La serie
$\sum_k (x_{n_{k+1}} - x_{n_k})$ converge absolutamente y, por tanto,
converge; sus sumas parciales son $x_{n_{k+1}} - x_{n_1}$, de modo que
$(x_{n_k})$ converge, y una sucesión de Cauchy con una subsucesión
convergente converge.
\end{solution}

\begin{solution}{exo:b3:complete:2}
(a) Supóngase $E$ \omterm{def:b3:complete:complete}{completo} con base
\omterm{def:b3:galois:algebraic}{algebraica} $(e_n)_{n\in\N}$ y sea
$F_n = \operatorname{Vect}(e_1, \dots, e_n)$: es cerrado (los
subespacios de dimensión finita son
\omterm{def:b3:complete:complete}{completos} y, por tanto, cerrados ---
segundo año) y de \omterm{def:b3:topology:interior}{interior} vacío: si
$B(x, r) \subseteq F_n$, tómese $v \notin F_n$; entonces
$x + \frac{r}{2\norm v}v \in B(x,r)
\setminus F_n$, absurdo. Pero $E = \bigcup F_n$ (todo vector es
combinación finita): contradice a Baire
(el~\cref{thm:b3:complete:baire}). El espacio $\R[X]$ tiene la base
numerable $(X^n)$, de modo que ninguna norma lo hace
\omterm{def:b3:complete:complete}{completo}.

(b) Fíjense $\delta > 0$ y una bola abierta no vacía $B$; buscaremos en
$B$ un punto de $\Omega_\delta = \{x : \operatorname{osc}_xf <
\delta\}$, donde $\operatorname{osc}_xf = \inf_{V \ni
x}\operatorname{diam} f(V)$. ($\Omega_\delta$ es abierto: si
$\operatorname{diam}f(V) < \delta$ para un abierto $V \ni x$, todo
$y \in V$ tiene oscilación $< \delta$.) Los conjuntos
\[
F_N = \bigl\{x : \abs{f_n(x) - f_m(x)} \leq \tfrac\delta3\
\ \forall\, n, m \geq N \bigr\}
\]
son cerrados (intersecciones de preimágenes de cerrados) y recubren $\R$
(la convergencia puntual hace $(f_n(x))$ de Cauchy). Aplicando Baire
dentro del \omterm{def:b3:complete:complete}{completo} $\bar B$: algún
$F_N \cap
\bar B$ contiene una bola $B' = B(x_0, \rho)$. Haciendo $m \to
\infty$: $\abs{f_N - f} \leq \frac\delta3$ sobre $B'$. Por la
\omterm{def:b3:topology:continuity}{continuidad} de $f_N$ en $x_0$,
redúzcase a $B'' \ni x_0$ donde
$\abs{f_N - f_N(x_0)} \leq \frac\delta3$; entonces, para $x \in B''$,
\[
\abs{f(x) - f(x_0)} \leq \abs{f - f_N}(x) + \abs{f_N(x) -
f_N(x_0)} + \abs{f_N - f}(x_0) \leq \delta:
\]
$\operatorname{diam} f(B'') \leq 2\delta$, luego $x_0 \in
\Omega_{3\delta} \cap B$. Así pues, cada $\Omega_\delta$ es abierto y
denso; $\bigcap_k\Omega_{1/k}$ --- el conjunto de los puntos de
\omterm{def:b3:topology:continuity}{continuidad} --- es denso por Baire.
\end{solution}

\begin{solution}{exo:b3:complete:3}
(a) $f(x) = \frac12(x + \frac ax)$ lleva $[\sqrt a, \infty)$ en sí mismo
(media aritmética--geométrica:
$f(x) \geq \sqrt{x \cdot \frac ax} = \sqrt a$) y allí
$f'(x) = \frac12(1 - \frac a{x^2}) \in [0, \frac12)$: es una contracción
$\frac12$-lipschitziana de un conjunto cerrado (y por tanto
\omterm{def:b3:complete:complete}{completo}). Punto fijo:
$x = f(x) \iff x^2 = a$: $x^* = \sqrt a$. Velocidad
(el~\cref{thm:b3:complete:banach}): $d(x_n, \sqrt a) \leq
2^{-n+1}\,d(x_1, x_0)$. Para $a = 2$, $x_0 = 2$: $d(x_1, x_0) =
\frac12$, de modo que $n = 40$ garantiza $2^{-40} < 10^{-12}$. (En
realidad el método de Newton converge cuadráticamente: bastan unas pocas
iteraciones; la estimación por contracción es pesimista, pero gratuita.)

(b) $f_m(x) = m + e\sin x$ es $e$-lipschitziana con $e < 1$ sobre el
\omterm{def:b3:complete:complete}{completo} $\R$: punto fijo único
$x(m)$. Para dos parámetros:
\[
\abs{x(m) - x(m')} = \abs{f_m(x(m)) - f_{m'}(x(m'))}
\leq \abs{m - m'} + e\abs{x(m) - x(m')},
\]
luego $\abs{x(m) - x(m')} \leq \frac{\abs{m - m'}}{1 - e}$: incluso
lipschitziana en $m$.
\end{solution}

\begin{solution}{exo:b3:complete:4}
(a) Sea $x^*$ el único punto fijo de $f^p$. Entonces
$f^p(f(x^*)) = f(f^p(x^*)) = f(x^*)$: $f(x^*)$ es un punto fijo de
$f^p$, luego $f(x^*) = x^*$. Y un punto fijo de $f$ lo es de $f^p$: la
unicidad se transfiere. Para $T$: por inducción, $\abs{T^pf(x)}
\leq \norm f_\infty x^p/p!$ (cada integración añade un factor
$\frac xk$), luego $\norm{T^p} \leq a^p/p!$. La aplicación $S(u) = g +
\lambda Tu$ cumple $S^p(u) - S^p(v) = \lambda^pT^p(u - v)$, de norma
$\leq \abs\lambda^pa^p/p!\,\norm{u - v} \to 0$: alguna $S^p$ es una
contracción, y $S$ tiene un único punto fijo: la ecuación de Volterra
$u = g + \lambda Tu$ es unívocamente
\omterm{def:b3:groups:derived}{resoluble} para \emph{todo} $\lambda$.

(b) $\varphi(x) = d(x, f(x))$ es
\omterm{def:b3:topology:continuity}{continua} sobre el
\omterm{def:b3:topology:compact}{compacto} $K$: alcanza su mínimo en
cierto $x_0$. Si $f(x_0) \neq x_0$, entonces
$\varphi(f(x_0)) = d(f(x_0), f^2(x_0)) < d(x_0, f(x_0)) =
\min\varphi$: absurdo. Unicidad: dos puntos fijos $x \ne y$ darían
$d(x,y) = d(f(x), f(y)) < d(x,y)$. Sin \omterm{def:b3:topology:compact}{compacidad}: $f(x) = x + \frac1x$
sobre $[1, \infty)$ cumple, para $x < y$,
$f(y) - f(x) = (y - x)\bigl(1 - \frac1{xy}\bigr) <
y - x$, y sin embargo $f(x) > x$ siempre: no hay punto fijo --- el
decrecimiento estricto de las distancias es más débil que un factor de
contracción uniforme.
\end{solution}

\begin{solution}{exo:b3:complete:5}
(a) El conjunto $A = \{g = h\}$ es la preimagen de la diagonal
$\Delta_Y$ por la \omterm{def:b3:topology:continuity}{aplicación
continua} $x \mapsto (g(x), h(x))$;
$\Delta_Y$ es cerrada porque $Y$ es
\omterm{def:b3:topology:hausdorff}{de Hausdorff} (para $y_1 \neq y_2$,
\omterm{def:b3:topology:topology}{entornos} abiertos disjuntos dan una
caja abierta alrededor de $(y_1,
y_2)$ disjunta de la diagonal, de modo que el complementario de
$\Delta_Y$ es abierto): así, $A$ es cerrado, contiene un
\omterm{def:b3:topology:interior}{conjunto denso} y es igual a $X$. Se
usó en la unicidad del~\cref{thm:b3:complete:extension} y, por tanto, en
la unicidad de los \omterm{thm:b3:complete:completion}{completados}.

(b) $f$, que es una isometría, es uniformemente
\omterm{def:b3:topology:continuity}{continua}: se extiende a
$F\colon X \to Y$ (el~\cref{thm:b3:complete:extension}) y sigue siendo
isométrica (la relación $d(F(x), F(x')) = d(x, x')$ vale sobre un
\omterm{def:b3:topology:interior}{conjunto denso} de pares y ambos
miembros son \omterm{def:b3:topology:continuity}{continuos}).
Análogamente, $f^{-1}\colon f(D) \to X$ se extiende a
$G \colon Y \to X$. La compuesta $G \circ F$ es
\omterm{def:b3:topology:continuity}{continua} y fija el denso $D$: es
$\mathrm{id}_X$ (apartado (a)); simétricamente $F \circ G =
\mathrm{id}_Y$. Luego $F$ es una biyección isométrica. Unicidad de los
\omterm{thm:b3:complete:completion}{completados}: aplíquese esto a
$D = X$, que está densamente en dos
\omterm{thm:b3:complete:completion}{completados}.
\end{solution}

\begin{solution}{exo:b3:complete:6}
$\{\sin(nx)\}$: puntualmente acotada por $1$; \emph{no}
\omterm{def:b3:complete:equicontinuous}{equicontinua}: en $x = 0$,
$\sin(n\cdot\frac{\pi}{2n}) = 1$ con $\frac\pi{2n} \to 0$, lo que viola
cualquier $\delta$ común para $\varepsilon = \frac12$. No relativamente
compacta (necesidad en Ascoli,
el~\cref{thm:b3:complete:ascoli}).

$\{x^n\}$: acotada; no
\omterm{def:b3:complete:equicontinuous}{equicontinua} en $1$: $1 - (1 -
\delta)^n \to 1$ cuando $n \to \infty$ para $\delta$ fijo. No
relativamente compacta ---
coherentemente, su límite puntual es discontinuo, de modo que ninguna
subsucesión converge uniformemente.

$\{\norm f_\infty \leq 1, \norm{f'}_\infty \leq 1\}$: la desigualdad del
valor medio hace la familia $1$-lipschitziana y, por tanto,
\omterm{def:b3:complete:equicontinuous}{equicontinua}; y es acotada:
relativamente compacta por Ascoli. (No
compacta: no es cerrada --- los
límites uniformes no tienen por qué ser $\mathcal C^1$; su
\omterm{def:b3:topology:interior}{clausura} son las funciones
$1$-lipschitzianas de norma $\leq 1$.)

$\{\operatorname{Lip}(f) \leq 1, f(0) = 0\}$:
\omterm{def:b3:complete:equicontinuous}{equicontinua}; puntualmente
acotada ($\abs{f(x)} \leq x \leq 1$); y cerrada por límites uniformes
(la desigualdad de Lipschitz y el valor en $0$ pasan al límite):
compacta.
\end{solution}

\begin{solution}{exo:b3:complete:7}
(a) Para $\norm f_\infty \leq 1$: $\abs{Tf(x)} \leq
\norm k_\infty$, y
\[
\abs{Tf(x) - Tf(x')} \leq \int_0^1\abs{k(x,y) -
k(x',y)}\,\dd y \leq \omega_k\bigl(\abs{x - x'}\bigr),
\]
donde $\omega_k$ es un módulo de
\omterm{def:b3:topology:continuity}{continuidad} uniforme de $k$ sobre
el cuadrado \omterm{def:b3:topology:compact}{compacto} (Heine): la
imagen de la bola unidad está uniformemente acotada y es
\omterm{def:b3:complete:equicontinuous}{equicontinua}, luego
relativamente compacta (Ascoli). Por
linealidad, todo acotado tiene imagen relativamente
compacta: $T$ es un operador
\omterm{def:b3:topology:compact}{compacto}.

(b) El enunciado sobre subsucesiones es la definición de \omterm{def:b3:topology:compact}{compacidad}
relativa aplicada a $\{Tf_n\}$. Si $T$ fuera biyectivo con inversa
\omterm{def:b3:topology:continuity}{continua}, entonces
$T(B(0,1)) \supseteq B(0,
\varepsilon)$ para cierto $\varepsilon > 0$ ($T^{-1}$
\omterm{def:b3:topology:continuity}{continua} en $0$); la bola cerrada
$\bar B(0,\varepsilon)$, subconjunto cerrado del
\omterm{def:b3:topology:compact}{compacto} $\overline{T(B(0,1))}$, sería
compacta --- imposible en el $\mathcal
C(\intcc01)$ de dimensión infinita por el teorema de Riesz (segundo
año).
\end{solution}

\begin{solution}{exo:b3:complete:8}
(a) Dini: véase el~\cref{lem:b3:complete:dini} --- el argumento de
recubrimiento. (b) Falso: la joroba móvil $f_n(x) = \max(0, 1 -
\abs{nx - 1})$ tiende a $0$ puntualmente (para $x > 0$, $f_n(x) =
0$ en cuanto $n > 2/x$; $f_n(0) = 0$) pero $\norm{f_n}_\infty = 1$. (c)
Falso: $f_n(x) = x^n$ decrece hacia la función discontinua
$\mathbf 1_{\{1\}}$; $\sup_{[0,1]}\abs{f_n - f} = \sup_{[0,1)}
x^n = 1$. (d) Falso: $x^n \downarrow 0$ puntualmente sobre el no
\omterm{def:b3:topology:compact}{compacto} $\intoo01$, con $\sup_{(0,1)}x^n = 1$. Cada hipótesis de Dini
es necesaria.
\end{solution}

\begin{solution}{exo:b3:complete:9}
(a) Por linealidad, $\int_0^1 fP = 0$ para todo polinomio $P$. Elíjanse
$P_n \to f$ uniformemente (el~\cref{cor:b3:complete:weierstrass}):
$\int_0^1 f^2 =
\lim\int_0^1 fP_n = 0$, y la función
\omterm{def:b3:topology:continuity}{continua} $f^2 \geq 0$ con integral
nula es idénticamente nula.

(b) Sobre $\intcc01$: los polinomios en $x^2$ forman un álgebra con
constantes que separa puntos ($x \mapsto x^2$ es inyectiva sobre
$\intcc01$): densa por Stone--Weierstrass. Sobre $\intcc{-1}1$: $x^2$
toma valores iguales en $\pm x$, y otro tanto le ocurre a todo polinomio
en $x^2$: un límite uniforme de tales polinomios es una función par. Si
funciones pares $g_n$ convergieran uniformemente a la identidad,
entonces $x = \lim g_n(x) = \lim g_n(-x) = -x$ para todo $x$: absurdo
--- no es densa. La hipótesis de separación falla en los pares
$\{x, -x\}$.

(c) No. Para $f$ en el álgebra $\mathcal A$ generada por las constantes
y $\eu^{\iu t}$ --- combinaciones lineales de $\eu^{\iu nt}$, $n \geq 0$
--- se tiene $\Lambda(f) =
\int_0^{2\pi}f(t)\,\eu^{\iu t}\,\dd t = 0$ (cada
$\int_0^{2\pi}\eu^{\iu(n+1)t}\dd t = 0$). $\Lambda$ es
\omterm{def:b3:topology:continuity}{continua} para $\norm\cdot_\infty$
($\abs{\Lambda(f)}\leq
2\pi\norm f_\infty$), de modo que $\Lambda$ se anula sobre
$\bar{\mathcal
A}$; pero $\Lambda(\eu^{-\iu t}) = 2\pi \neq 0$: $\eu^{-\iu t}
\notin \bar{\mathcal A}$. (Stone--Weierstrass no se aplica: $\mathcal A$
no es estable por conjugación --- y la obstrucción es precisamente la
que la teoría de funciones holomorfas sistematizará en
el~\cref{ch:b3:holomorphic}.)
\end{solution}

\begin{solution}{exo:b3:complete:10}
$F_n = \bigcap_{f \in \mathcal F}\{x : \abs{f(x)} \leq n\}$ es una
intersección de cerrados: es cerrado. La acotación puntual da
$X = \bigcup_n F_n$. Baire (el~\cref{thm:b3:complete:baire}) proporciona
$n_0$ con $U = \mathring F_{n_0} \neq \varnothing$: sobre $U$,
$\abs f \leq n_0$ para todo $f \in \mathcal F$ simultáneamente.
\end{solution}

\begin{solution}{exo:b3:complete:11}
(a) Una sucesión de Cauchy para $\norm\cdot_{\mathcal C^1}$ es
uniformemente de Cauchy junto con sus derivadas: $f_n \to f$ y
$f_n' \to g$ uniformemente, con $f, g$
\omterm{def:b3:topology:continuity}{continua}. Pasando al límite (la
convergencia uniforme lo permite bajo la integral) en
$f_n(x) = f_n(0) + \int_0^xf_n'(t)\,\dd t$ se obtiene
$f(x) = f(0) + \int_0^xg$: $f$ es $\mathcal C^1$ con $f' = g$, y
$\norm{f_n - f}_{\mathcal C^1} \to 0$.
\omterm{def:b3:complete:complete}{Completo}.

(b) $h(x) = \abs{x - \frac12}$ es límite uniforme de funciones
$\mathcal C^1$, por ejemplo $h_n(x) = \sqrt{(x -
\frac12)^2 + \frac1n}$ ($\abs{h_n - h} \leq \frac1{\sqrt n}$ por la cota
de la cantidad conjugada), y sin embargo $h \notin E$: la norma $\sup$
sobre $E$ no es \omterm{def:b3:complete:complete}{completa} --- su
\omterm{thm:b3:complete:completion}{completado} es
$\mathcal C(\intcc01)$.

(c) $f_n(x) = \sin(2\pi nx)$ cumple $\norm{f_n}_\infty = 1$ y
$\norm{f_n'}_\infty = 2\pi n \to \infty$: no existe ningún $C$.
(Conceptualmente: si la derivación estuviera acotada para la norma del
supremo, las dos normas de (a)--(b) serían equivalentes, lo que haría
$(E, \norm\cdot_\infty)$ \omterm{def:b3:complete:complete}{completo} ---
en contra de (b). Esta falta de acotación es exactamente la razón de que
las funciones \omterm{def:b3:topology:continuity}{continuas} genéricas
puedan no ser derivables en ningún punto,
\cref{thm:b3:complete:nowherediff}.)
\end{solution}

\begin{solution}{exo:b3:complete:12}
Fíjese $\varepsilon > 0$. Cada $F_N$ es una intersección sobre $n
\geq N$ de preimágenes del cerrado $\intcc{-\varepsilon}
\varepsilon$ por la \omterm{def:b3:topology:continuity}{aplicación
continua} $x \mapsto f(nx)$: es
cerrado. La hipótesis dice que todo $x > 0$ está en algún $F_N$. Por
Baire aplicado dentro del \omterm{def:b3:complete:complete}{completo}
$\intcc ab$ (cualquier $0 < a <
b$), algún $F_N$ es denso en un subintervalo; y, siendo cerrado,
contiene un intervalo $\intcc{a'}{b'}$ con $0 < a' < b'$. Entonces, para
todo $n \geq \max(N, \frac{a'}{b' - a'})$, los intervalos
$\intcc{na'}{nb'}$ y $\intcc{(n+1)a'}{(n+1)b'}$ se solapan
($nb' \geq (n+1)a'$), luego
\[
\bigcup_{n \geq n_0}\intcc{na'}{nb'} \supseteq
\intco{n_0a'}{+\infty},
\]
y todo $t \geq n_0a'$ es $t = nx$ con $n \geq N$, $x \in
\intcc{a'}{b'} \subseteq F_N$: $\abs{f(t)} \leq \varepsilon$. Por tanto,
$\limsup_{t\to\infty}\abs f \leq \varepsilon$ para todo $\varepsilon$:
$f \to 0$. La hipótesis \omterm{def:b3:complete:complete}{completa} es necesaria porque los $F_N$ han de
\emph{recubrir} todo un intervalo de valores de $x$ --- solo con $x$
racionales, la unión de los $F_N$ sería numerable y Baire no daría nada;
de hecho, hay contraejemplos
\omterm{def:b3:topology:continuity}{continuos} que se anulan a lo largo
de todas las semirrectas racionales pero no en el infinito.
\end{solution}

\begin{solution}{pb:b3:complete:1}
\textbf{1.} Inducción sobre $k$: si $\norm{\varphi_n(t_k) - x_0}
\leq M(t_k - t_0) \leq MT \leq b$, el punto $(t_k,
\varphi_n(t_k))$ está en $R$, de modo que la pendiente $f(t_k,
\varphi_n(t_k))$ está definida, de norma $\leq M$; entonces, para $t \in
[t_k, t_{k+1}]$, $\norm{\varphi_n(t) - x_0} \leq
\norm{\varphi_n(t_k) - x_0} + M(t - t_k) \leq M(t - t_0) \leq
MT \leq b$.

\textbf{2.} Cada trozo afín tiene pendiente de norma $\leq M$; y una
función afín a trozos con pendientes acotadas por $M$ es
$M$-lipschitziana (encadénese a través de los puntos de la malla).

\textbf{3.} La familia $(\varphi_n)$ está puntualmente acotada (valores
en $\bar B(x_0, b)$) y es
\omterm{def:b3:complete:equicontinuous}{equicontinua} (constante de
Lipschitz común $M$): Ascoli (el~\cref{thm:b3:complete:ascoli}) extrae
$\varphi_{n_j} \to
\varphi$ uniformemente sobre $[t_0, t_0 + T]$. Las cotas pasan al
límite: $\varphi$ es $M$-lipschitziana y $\varphi(t_0) = x_0$.

\textbf{4.} $R$ es \omterm{def:b3:topology:compact}{compacto} y $f$ es
\omterm{def:b3:topology:continuity}{continua}: es uniformemente
\omterm{def:b3:topology:continuity}{continua} (Heine,
el~\cref{cor:b3:topology:heineborel}); sea $\delta(\varepsilon)$ un
módulo. Para $t \in (t_k,
t_{k+1})$ fuera de la malla: $\varphi_n'(t) = f(t_k, \varphi_n(t_k))$, y
los dos puntos de evaluación de $f$ difieren en $\abs{t - t_k} \leq T/n$
en tiempo y en $\norm{\varphi_n(t) - \varphi_n(t_k)} \leq MT/n$ en
espacio. Para $n \geq n_0$ con $(1 + M)T/n_0 < \delta(\varepsilon)$:
$\norm{\Delta_n(t)} \leq \varepsilon$.

\textbf{5.} Sobre cada $[t_k, t_{k+1}]$, $\varphi_n$ es afín, luego
$\varphi_n(t_{k+1}) - \varphi_n(t_k) = \int_{t_k}^{t_{k+1}}
\varphi_n'(s)\dd s$, siendo la derivada la pendiente constante; sumando
sobre los trozos (y cortando el último en $t$):
$\varphi_n(t) = x_0 + \int_{t_0}^t\varphi_n'(s)\,\dd s$.
Escribiendo $\varphi_n' = f(s, \varphi_n(s)) + \Delta_n(s)$ (integrandos
\omterm{def:b3:topology:continuity}{continuos} a trozos, con finitos
saltos) se obtiene la fórmula.

\textbf{6.} Dado $\varepsilon$: para $j$ grande,
$\norm{\varphi_{n_j} - \varphi}_\infty < \delta(\varepsilon)$,
luego $\norm{f(s, \varphi_{n_j}(s)) - f(s, \varphi(s))} \leq
\varepsilon$ para todo $s$: convergencia uniforme de los integrandos, y
$\int_{t_0}^t f(s,\varphi_{n_j}(s))\dd s \to
\int_{t_0}^tf(s, \varphi(s))\dd s$ uniformemente en $t$. Además,
$\norm{\int_{t_0}^t\Delta_{n_j}} \leq T\sup\norm{\Delta_{n_j}}
\to 0$ (pregunta 4). Pasando al límite en la identidad de la pregunta 5:
$\varphi(t) = x_0 + \int_{t_0}^tf(s,
\varphi(s))\,\dd s$.

\textbf{7.} El integrando $s \mapsto f(s, \varphi(s))$ es
\omterm{def:b3:topology:continuity}{continuo}, de modo que el miembro
derecho es $\mathcal C^1$ en $t$ con derivada $f(t, \varphi(t))$:
$\varphi$ resuelve el problema de Cauchy sobre $[t_0, t_0 + T]$. Para la
mitad izquierda, póngase $g(t, x) =
-f(2t_0 - t, x)$, \omterm{def:b3:topology:continuity}{continua} sobre el
rectángulo reflejado con la misma cota $M$; una solución $\psi$ de
$y' = g(t,y)$, $y(t_0) = x_0$ sobre $[t_0, t_0 + T]$ da $\varphi(t) =
\psi(2t_0 - t)$, que resuelve la ecuación original sobre $[t_0 - T,
t_0]$; y las dos mitades se pegan en una solución $\mathcal C^1$ (ambas
derivadas laterales en $t_0$ valen $f(t_0, x_0)$). --- \emph{Teorema de
Peano}: una $f$ \omterm{def:b3:topology:continuity}{continua} admite una
solución local con cualquier condición inicial.

\textbf{8.} La \omterm{def:b3:topology:continuity}{continuidad} es
clara. La condición de Lipschitz cerca de $0$ falla:
$\abs{f(h) - f(0)} = 2\sqrt h$, y $2\sqrt h \leq L h$ es falso para
$h < 4/L^2$.

\textbf{9.} Para $t \leq c$: $x_c \equiv 0$ es solución. Para $t \geq
c$: $x_c'(t) = 2(t - c) = 2\sqrt{(t-c)^2} = 2\sqrt{\abs{x_c}}$. En
$t = c$, ambas derivadas laterales valen $0$: $x_c$ es $\mathcal C^1$ y
resuelve globalmente, con $x_c(0) = 0$ para todo $c \geq 0$ --- junto
con $x \equiv 0$, un \omterm{def:b3:topology:continuity}{continuo} de soluciones que pasan por el origen.

\textbf{10.} La iteración de Picard plantea $\Phi(x)(t) = x_0 +
\int_0^t f(x(s))\dd s$ y necesita $\norm{\Phi(x) - \Phi(y)}
\leq k\norm{x - y}$ con $k < 1$ sobre una bola adecuada --- lo cual se
sigue de una cota de Lipschitz sobre $f$, trasladada bajo la integral.
Aquí $\sqrt\cdot$ no admite ninguna cota de Lipschitz cerca de $0$, y
ninguna elección de intervalo o de bola lo repara. Y, en efecto, ninguna
demostración de unicidad podría funcionar: la unicidad es \emph{falsa}
(pregunta 9).

\textbf{11.} En $c_0$, los vectores unidad $e_n = (0, \dots, 0, 1,
0, \dots)$ cumplen $\norm{e_n - e_m}_\infty = 1$ para $n \neq
m$: ninguna subsucesión es de Cauchy, de modo que la bola unidad cerrada
no es compacta. El paso que se rompe
es la extracción (pregunta 3): Ascoli para $\mathcal C([t_0, t_0+T], E)$
exige que los valores vivan en un espacio donde los acotados sean
relativamente \omterm{def:b3:topology:compact}{compactos} --- cierto en
$\R^d$ (Bolzano--Weierstrass), falso en $c_0$; la extracción puntual
deja de estar disponible (de hecho, el ejemplo de Dieudonné no tiene
ninguna solución local).

\textbf{12.} \emph{Peano}: si $f$ es
\omterm{def:b3:topology:continuity}{continua} en un
\omterm{def:b3:topology:topology}{entorno} de $(t_0, x_0)$ en
$\R\times\R^d$, $\Rightarrow$ existe una solución $\mathcal C^1$ de
$x' = f(t,x)$, $x(t_0) = x_0$, sobre cierto $[t_0 - T, t_0 + T]$.
\emph{Cauchy--Lipschitz} (el~\cref{ch:b3:ode}): si además $f$ es
localmente lipschitziana en la variable $x$, la solución es \emph{única}
(dos cualesquiera coinciden en su intervalo común) --- existencia y
unicidad. El par $(x' = 2\sqrt{\abs x},\ x(0) = 0)$ separa los dos
teoremas.

\textbf{13.} $\omega(r) = Lr$: $\int_0^1\frac{\dd r}{Lr} =
+\infty$: sirve. $\omega(r) = r\log\frac1r$ (cerca de $0$):
$\int\frac{\dd r}{r\log(1/r)} = \bigl[-\log\log\frac1r\bigr]
\to +\infty$ cuando $r \to 0$: sirve --- y sin embargo $\omega(r)/r =
\log\frac1r \to \infty$: no es lipschitziana. $\omega(r) = 2\sqrt
r$: $\int_0^1\frac{\dd r}{2\sqrt r} = \bigl[\sqrt r\bigr]_0^1
= 1 < \infty$: no sirve.

\textbf{14.} Réstense las dos formas integrales:
\[
x_1(t) - x_2(t) = x_1(s) - x_2(s) + \int_s^t\bigl(f(v,
x_1(v)) - f(v, x_2(v))\bigr)\dd v,
\]
tómense normas y úsese el módulo de Osgood: $\delta(t) \leq
\delta(s) + \int_s^t\omega(\delta(v))\,\dd v$.

\textbf{15.} $\tau$ está bien definida ($\delta(t_0) = 0$) con
$\delta(\tau) = 0$ (\omterm{def:b3:topology:continuity}{continuidad}) y
$\delta > 0$ sobre $\intoc\tau{t_1}$. Fíjese $s \in \intoo\tau{t_1}$.
Entonces $u$ es $\mathcal C^1$, $u(s) = \delta(s) > 0$, $u \geq \delta$
sobre $[s, t_1]$ (pregunta 14), y $u' = \omega(\delta) \leq
\omega(u)$ ($\omega$ no decreciente, $u > 0$). Divídase e intégrese:
\[
\int_{\delta(s)}^{u(t_1)}\frac{\dd r}{\omega(r)}
= \int_s^{t_1}\frac{u'(v)}{\omega(u(v))}\,\dd v \leq t_1 - s
\leq t_1 - \tau .
\]
Aunque $u$ depende de $s$, la cota $u(t_1) \geq
\delta(t_1)$ vale para todo $s$, de modo que el miembro izquierdo es al
menos $\int_{\delta(s)}^{\delta(t_1)}\frac{\dd r}{\omega(r)}$, que
tiende a $+\infty$ cuando $s \downarrow \tau$ (entonces
$\delta(s) \to \delta(\tau) = 0$, y la integral diverge en $0$): el
miembro derecho, acotado, queda contradicho. Por tanto,
$\delta \equiv 0$: unicidad.

\textbf{16.} Lipschitz es $\omega = Lr$: se recupera la unicidad. Para
$x' = x\log\frac1{\abs x}$: el miembro derecho cumple el módulo de
Osgood $\omega(r) = r\log\frac1r$ cerca de $0$ (desigualdad del valor
medio sobre $x \mapsto x\log\frac1x$, cuya derivada $\log\frac1x - 1$ no
está acotada --- Lipschitz falla, Osgood se cumple): soluciones únicas;
obsérvese que $x \equiv 0$ es una de ellas, de modo que ninguna otra
solución puede tocar $0$. Para $\omega(r) = 2\sqrt r$:
$\int_0^{h^2}\frac{\dd r}{2\sqrt r} = h$ --- el «presupuesto de Osgood»
para subir de $0$ a la altura $h^2$ es exactamente el tiempo $h$ y, en
efecto, $x_c(c + h) = h^2$: la solución emplea el tiempo $h$ en hacer
precisamente lo que la integral convergente permite. La divergencia de
la integral es la imposibilidad de abandonar $0$ en tiempo finito; la
convergencia es la vía de escape.

\textbf{17.} $G(t) = \int_{t_0}^t(Le(s) + \eta(s))\dd s$ es
$\mathcal C^1$ con $G' = Le + \eta \leq LG + \sup\eta$ (hipótesis
$e \leq G$). Entonces $\bigl(\eu^{-Lt}(G +
\tfrac{\sup\eta}L)\bigr)' = \eu^{-Lt}(G' - LG - \sup\eta)
\leq 0$: el corchete decrece, luego $G(t) + \frac{\sup\eta}L
\leq \eu^{L(t - t_0)}\bigl(G(t_0) + \frac{\sup\eta}L\bigr) =
\eu^{L(t-t_0)}\frac{\sup\eta}L$, es decir, $e(t) \leq G(t) \leq
\frac{\sup\eta}L(\eu^{L(t-t_0)} - 1)$.

\textbf{18.} Defecto cuantitativo: sobre $(t_k, t_{k+1})$,
$\Delta_n(t) = f(t_k, \varphi_n(t_k)) - f(t, \varphi_n(t))$
con $\abs{t - t_k} \leq T/n$ y $\norm{\varphi_n(t) -
\varphi_n(t_k)} \leq MT/n$, luego $\norm{\Delta_n} \leq L'T/n +
LMT/n$. Restando las identidades integrales de $\varphi_n$ (pregunta 5)
y $\varphi$ y usando la cota de Lipschitz:
\[
\norm{\varphi_n(t) - \varphi(t)} \leq
\int_{t_0}^t L\,\norm{\varphi_n - \varphi}(s)\,\dd s +
\int_{t_0}^t\norm{\Delta_n(s)}\,\dd s,
\]
y la pregunta 17 con $\eta = \norm{\Delta_n}$ da la cota $O(1/n)$
enunciada, siendo $\varphi$ única por Cauchy--Lipschitz (o por Osgood).
Incluso sin velocidades: toda subsucesión de las $(\varphi_n)$,
equiacotadas y equilipschitzianas, tiene una subsubsucesión convergente
(Ascoli y Parte II) hacia \emph{alguna} solución, que la unicidad obliga
a ser $\varphi$: y una sucesión en la que toda subsucesión tiene una
subsubsucesión con el mismo límite converge.

\textbf{19.} De $\varphi_n(t_k) = 0$: la pendiente $2\sqrt0 =
0$ da $\varphi_n(t_{k+1}) = 0$; por inducción, $\varphi_n
\equiv 0$, que converge a la solución nula. De $\varepsilon > 0$: sobre
$[\varepsilon', \infty)$ con $\varepsilon' < \varepsilon$, la función
$2\sqrt x$ es lipschitziana, de modo que se aplica la pregunta 18 y
Euler converge a la única solución que pasa por $(0, \varepsilon)$, a
saber $x(t) =
(t + \sqrt\varepsilon)^2$ (compruébese: $x' = 2(t + \sqrt
\varepsilon) = 2\sqrt x$). Cuando $\varepsilon \to 0$, $(t +
\sqrt\varepsilon)^2 \to t^2$ uniformemente sobre $[0,1]$: el límite
doble aterriza en $x_0(t) = t^2$, y no en $0$. Una perturbación
arbitrariamente pequeña del dato inicial redirige el esquema de una
solución a otra: la no unicidad leída como inestabilidad numérica.

\textbf{20.} No vacío: Parte II. Toda solución cumple
$\norm{x'} = \norm{f(t, x)} \leq M$: $\mathcal S$ es uniformemente
$M$-lipschitziano y uniformemente acotado (valores en $\bar B(x_0, b)$).
Cerrado: si $x_n \in \mathcal S \to x$ uniformemente, pásese al límite
en $x_n(t) = x_0 +
\int_{t_0}^tf(s, x_n(s))\dd s$ (los integrandos convergen uniformemente
por la \omterm{def:b3:topology:continuity}{continuidad} uniforme de $f$
sobre el \omterm{def:b3:topology:compact}{compacto} $R$):
$x \in \mathcal S$. Ascoli: $\mathcal S$ es un subconjunto cerrado,
acotado y equicontinuo de
$\mathcal C([t_0, t_0+T],
\R^d)$: \omterm{def:b3:topology:compact}{compacto}.

\textbf{21.} Toda solución es no decreciente ($x' = 2\sqrt{\abs
x} \geq 0$) con $x(0) = 0$, luego $x \geq 0$. Sea $c =
\sup\{t \in [0,1] : x(t) = 0\}$ (posiblemente $c = \infty$ si
$x \equiv 0$, en cuyo caso $x = x_\infty$). Para $t > c$: $x
> 0$ (monotonía y definición de $c$), y allí $(\sqrt
x)' = \frac{x'}{2\sqrt x} = 1$, luego $\sqrt{x(t)} = t - c$
(\omterm{def:b3:topology:continuity}{continuidad} en $c$): $x = x_c$.
\omterm{def:b3:topology:continuity}{Continuidad} de $c \mapsto
x_c$: para $c, c' \in [0, 1]$, $\sup_{[0,1]}\abs{(t - c)_+^2 -
(t - c')_+^2} \leq 2\abs{c - c'}$ (la aplicación $c \mapsto
(t-c)_+^2$ es $2$-lipschitziana uniformemente en $t \in [0,1]$) y, como
$x_c \equiv 0$ sobre $[0,1]$ para todo $c \geq 1$, la familia se reduce
a $\mathcal S = \{x_c : c \in [0,1]\}$ (con $x_1 = 0 = x_\infty$). Así
pues, $\mathcal S$ es la imagen del
\omterm{def:b3:topology:compact}{compacto}
\omterm{def:b3:topology:connected}{conexo} $[0,1]$ por la
\omterm{def:b3:topology:continuity}{aplicación continua} $c
\mapsto x_c$: \omterm{def:b3:topology:compact}{compacto} y
\omterm{def:b3:topology:connected}{conexo}. El embudo de soluciones es
un segmento \omterm{def:b3:topology:continuity}{continuo} que va de
$t^2$ (escape inmediato) hasta $0$ (reposo eterno).

\textbf{22.} La evaluación $\operatorname{ev}_t\colon
\mathcal C([0,1]) \to \R$, $x \mapsto x(t)$, es
\omterm{def:b3:topology:continuity}{continua}
($\abs{x(t) - y(t)} \leq \norm{x - y}_\infty$), de modo que $\mathcal
S(t) = \operatorname{ev}_t(\mathcal S)$ es imagen
\omterm{def:b3:topology:continuity}{continua} de un
\omterm{def:b3:topology:compact}{compacto}: es
\omterm{def:b3:topology:compact}{compacto} --- y de un
\omterm{def:b3:topology:connected}{conexo}:
\omterm{def:b3:topology:connected}{conexo}. Para el ejemplo:
$x_c(t) = (t - c)_+^2$ barre, cuando $c$ recorre $[0, 1]$, todos los
valores desde $t^2$ (en $c = 0$) hasta $0$ (en $c \geq t$), de manera
\omterm{def:b3:topology:continuity}{continua}: $\mathcal S(t) =
\intcc0{t^2}$. En cada instante, la sección del embudo es un segmento
\omterm{def:b3:complete:complete}{completo}: entre el reposo y el escape máximo, todo compromiso lo realiza
alguna solución.

\textbf{23.} Restando la forma integral $x(t) = x_0 +
\int_{t_0}^tf(s, x(s))\dd s$ y su análoga para $y$, y poniendo
$e(t) = \norm{x(t) - y(t)}$, $d = \norm{x_0 - y_0}$:
\[
e(t) \leq d + \int_{t_0}^tL\,e(s)\dd s = G(t).
\]
Entonces $G(t_0) = d$ y $G' = Le \leq LG$, luego $\bigl(\eu^{-L(t
- t_0)}G\bigr)' \leq 0$ y $G(t) \leq d\,\eu^{L(t - t_0)}$; de donde
$e(t) \leq d\,\eu^{L(t-t_0)}$. Optimalidad: para $x' =
Lx$, las soluciones que pasan por $x_0$ y $y_0$ son $x_0\eu^{L(t -
t_0)}$ y $y_0\eu^{L(t-t_0)}$, cuya distancia es exactamente
$d\,\eu^{L(t-t_0)}$. Con $d = 0$, $e \equiv 0$: la unicidad de
Cauchy--Lipschitz, redemostrada en dos líneas. Y para
$t \in [t_0, t_0 + T]$ fijo, $\norm{x(t; x_0) - x(t;
y_0)} \leq \eu^{LT}\norm{x_0 - y_0}$: el flujo es lipschitziano en la
condición inicial --- dependencia determinista, a un precio exponencial
controlado.

\textbf{24.} Sobre $\intoc01$, $\Omega$ está bien definida (la integral
converge en $0$ por hipótesis), es $\mathcal C^1$ con
$\Omega' = \frac1\omega > 0$: una biyección creciente sobre
$\intoc0{\Omega(1)}$, con $\Omega(x) \to 0$ cuando $x \to 0^+$. Su
inversa $g \colon \intoc0{\Omega(1)} \to \intoc01$ es $\mathcal C^1$ con
\[
g'(t) = \frac1{\Omega'(g(t))} = \omega\bigl(g(t)\bigr) > 0,
\qquad g(t) \xrightarrow[t \to 0^+]{} 0 .
\]
Extiéndase $g$ por $0$ sobre $t \leq 0$: la
\omterm{def:b3:topology:continuity}{continuidad} es clara y, en $t = 0$,
para $t > 0$,
\[
\frac{g(t)}t = \frac1t\int_0^tg'(s)\dd s
= \frac1t\int_0^t\omega\bigl(g(s)\bigr)\dd s
\leq \omega\bigl(g(t)\bigr) \xrightarrow[t\to0^+]{} 0
\]
($\omega$ no decreciente, $g$ creciente, $\omega(0^+) = 0$):
$g'(0) = 0 = f(g(0))$, y $g' = f(g)$ se cumple a ambos lados de $0$. Así
pues, $x \equiv 0$ y $g$ son dos soluciones $\mathcal
C^1$ distintas que pasan por $(0, 0)$: cuando $\int_0\frac{\dd r}
{\omega(r)}$ converge, la unicidad falla --- la divergencia de la
pregunta 15 es exactamente la frontera. Para $\omega(r) = 2\sqrt
r$: $\Omega(x) = \int_0^x\frac{\dd r}{2\sqrt r} = \sqrt x$,
$g(t) = t^2$, y las traslaciones temporales dan toda la familia $x_c$ de
la Parte III.

\textbf{25.} Con paso $h = \frac1n$: $\varphi_n(t_{k+1}) =
\varphi_n(t_k) + h\,\varphi_n(t_k) = (1 + h)\varphi_n(t_k)$, luego
$\varphi_n(1) = (1 + \frac1n)^n$ tras $n$ pasos. Desarrollo:
\[
n\log\Bigl(1 + \frac1n\Bigr)
= n\Bigl(\frac1n - \frac1{2n^2} + O\Bigl(\frac1{n^3}\Bigr)
\Bigr) = 1 - \frac1{2n} + O\Bigl(\frac1{n^2}\Bigr),
\]
y, exponenciando, $(1 + \frac1n)^n = \eu\,\eu^{-1/(2n) +
O(n^{-2})} = \eu\bigl(1 - \frac1{2n} + O(n^{-2})\bigr)$. El error en
$t = 1$ es, por tanto, $\eu - \varphi_n(1) =
\frac{\eu}{2n} + O(n^{-2})$: el $O(\frac1n)$ de la pregunta 18, aquí con
su constante exacta $\frac\eu2$. Numéricamente, $n =
10$: $1.1^{10} = 2.5937424601$ ($1.1^2 = 1.21$, $1.1^4 =
1.4641$, $1.1^8 = 2.14358881$, por $1.21$), y $\eu -
2.59374 \approx 0.12454$, frente a la predicción asintótica
$\frac{\eu}{20} \approx 0.13591$: concordancia hasta la corrección
$O(n^{-2})$, cuyo término principal rebaja aquí la predicción hacia el
valor observado.
\end{solution}
